% =============================================================================
% SECTION 4 — EXPERIMENTS (revised for paper_v2)
% Tables are unchanged; prose tightened, ablation block (B) placeholders
% filled from the table, mask-granularity (block C) prose added.
% =============================================================================

\section{Experiments}
\label{sec:exp}
\vspace{-0.5ex}

\subsection{Setup}
\label{sec:exp:setup}

\textbf{Benchmarks.} We evaluate on three groups.
\emph{Coding-agent benchmarks}: SWE-Bench-Verified ($500$) and
SWE-Bench-Lite ($300$)~\citep{swebench}, our primary in-domain target.
\emph{Non-agent coding benchmarks}: LiveCodeBench~\citep{livecodebench},
OJBench~\citep{ojbench}, and FullStackBench-EN~\citep{fullstackbench};
these probe pure code generation and serve as a regression check.
\emph{Tool-use and OOD agent benchmarks}: Terminal-Bench~\citep{terminalbench},
$\tau$-bench~\citep{taubench}, and BFCL~\citep{bfcl}; the latter two
contain no Python code-editing trajectories and test whether the
function-call inductive bias transfers across task families.

\textbf{Mid-training and post-training pipeline.}
We mid-train Qwen2.5-Coder-Instruct (7B/14B)~\citep{qwencoder} and
Qwen3-8B~\citep{qwen3} (base, not Instruct) on the selected FIM corpus
using the standard FIM loss on the middle span
(rationale plus body), packed to the model's native context length and
using its native FIM sentinel tokens. We then apply existing agentic
post-training pipelines without modification:
R2E-Gym~\citep{r2egym} or SWE-Smith~\citep{swesmith} for the
Qwen2.5-Coder runs, and SWE-Lego~\citep{swelego} for Qwen3-8B. For
Qwen3-8B we train SWE-Lego for $2$ epochs rather than the official $4$
to prevent overfitting in our setup. Hyperparameters and token budgets
are in Appendix~\ref{app:hp}.

\textbf{Evaluation protocol.}
All numbers are means over three independent evaluation seeds on the
final checkpoint of each pipeline, with ``\%'' omitted in table cells
and std bands in parentheses. \emph{Baseline} is base $+$ post-training;
\emph{ours} is base $+$ FIM mid-training $+$ identical post-training.
We do not evaluate mid-training-only checkpoints because FIM-only models
have degraded instruction-following and cannot be compared fairly with
instruction-tuned baselines---every reported gain therefore
\emph{survives} subsequent post-training. The agent harness is fixed by
the post-training pipeline: R2E-Gym uses an OpenHands fork~\citep{openhands,
r2egym}, SWE-Smith uses SWE-agent~\citep{sweagent}, and the Qwen3-8B
runs use OpenHands directly (SWE-Lego data exceeds the 32K context of
the other scaffolds).

\textbf{Models and baselines.}
For Qwen2.5-Coder-Instruct at 7B/14B~\citep{qwencoder} we report
(i) the instruction-tuned base, (ii) post-training with R2E-Gym
reproduced under our setup, (iii) the official R2E-Gym numbers (in grey,
where available), (iv) our recipe of FIM mid-training followed by
R2E-Gym post-training, and at 7B additionally the analogous
SWE-Smith~\citep{swesmith} variants. Qwen3-8B uses the
SWE-Lego~\citep{swelego} pipeline; the cross-base-model comparison thus
varies the post-training pipeline simultaneously, a confound we discuss
in Section~\ref{sec:exp:main}.

% =============================================================================
% 4.2 Main Results
% =============================================================================
\vspace{-0.5ex}
\subsection{Main Results on SWE Agent Benchmarks}
\label{sec:exp:main}

\begin{table}[t]
\centering
\caption{Main results on coding agent benchmarks. All numbers are
percentages, means over three independent seeds with std bands;
\emph{Average} is the unweighted mean of Verified and Lite. Bolded rows
are our recipe; rows tagged \emph{officially reported} are quoted from
the corresponding original publications. SWE-Lego is post-trained for
$2$ epochs (vs.\ $4$ in the official release) to prevent overfitting.}
\label{tab:main}
\setlength{\tabcolsep}{6pt}
\begin{tabular}{lccc}
\toprule
Setting & SWE-Bench-Verified & SWE-Bench-Lite & Average \\
\midrule
\multicolumn{4}{l}{\textit{Qwen2.5-Coder-7B-Instruct}} \\
\quad --- (no agentic training)                       & 1.8\,($\pm$1.3)            & 1.0\,($\pm$1.0)            & 1.40 \\
\quad + R2E-Gym~\citep{r2egym} (reproduced)           & 15.00\,($\pm$1.5)          & 11.33\,($\pm$1.2)          & 13.17 \\
\quad + R2E-Gym (officially reported)                 & 19.00\,($\pm$1.0)          & 11.00\,($\pm$0.8)          & 15.00 \\
\quad \textbf{+ FIM-Midtrain + R2E-Gym}               & \textbf{17.80}\,($\pm$1.4) & \textbf{15.00}\,($\pm$1.1) & \textbf{16.40} \\
\rowcolor{blue!8}
\quad\quad $\Delta$ (ours vs.\ reproduced)            & $+2.80$                    & $+3.67$                    & $+3.24$ \\
\cmidrule(lr){1-4}
\quad + SWE-Smith~\citep{swesmith} (reproduced)       & 12.30\,($\pm$1.2)          & 14.20\,($\pm$1.4)          & 13.25 \\
\quad + SWE-Smith (officially reported)               & 15.20                       & 11.70                       & 13.45 \\
\quad \textbf{+ FIM-Midtrain + SWE-Smith}             & \textbf{17.60}\,($\pm$1.3) & \textbf{14.70}\,($\pm$1.0) & \textbf{16.15} \\
\rowcolor{blue!8}
\quad\quad $\Delta$ (ours vs.\ reproduced)            & $+5.30$                    & $+0.50$                    & $+2.90$ \\
\midrule
\multicolumn{4}{l}{\textit{Qwen2.5-Coder-14B-Instruct}} \\
\quad --- (no agentic training)                       & 4.0\,($\pm$1.6)            & 2.7\,($\pm$1.0)            & 3.35 \\
\quad + R2E-Gym (reproduced)                          & 26.20\,($\pm$1.4)          & 18.00\,($\pm$1.1)          & 22.10 \\
\quad + R2E-Gym (officially reported)                 & 26.80\,($\pm$1.4)          & 20.67\,($\pm$0.7)          & 23.74 \\
\quad \textbf{+ FIM-Midtrain + R2E-Gym}               & \textbf{29.20}\,($\pm$1.5) & \textbf{22.00}\,($\pm$1.2) & \textbf{25.60} \\
\rowcolor{blue!8}
\quad\quad $\Delta$ (ours vs.\ reproduced)            & $+3.00$                    & $+4.00$                    & $+3.50$ \\
\midrule
\multicolumn{4}{l}{\textit{Qwen3-8B}} \\
\quad --- (no agentic training)                       & 7.60\,($\pm$1.2)           & 5.80\,($\pm$0.9)           & 6.70 \\
\quad + SWE-Lego~\citep{swelego} (reproduced)         & 31.80\,($\pm$1.0)          & 27.30\,($\pm$1.1)          & 29.55 \\
\quad \textbf{+ FIM-Midtrain + SWE-Lego}              & \textbf{35.00}\,($\pm$1.5) & \textbf{32.70}\,($\pm$1.3) & \textbf{33.85} \\
\rowcolor{blue!8}
\quad\quad $\Delta$ (ours vs.\ reproduced)            & $+3.20$                    & $+5.40$                    & $+4.30$ \\
\bottomrule
\end{tabular}
\end{table}

Table~\ref{tab:main} summarizes in-domain agent results.

\textbf{Consistent gains on the Qwen2.5-Coder series.}
Holding the post-training pipeline fixed at R2E-Gym, FIM mid-training
improves SWE-Bench-Verified by $+2.80$ on 7B-Instruct and $+3.00$ on
14B-Instruct, with matching directional gains on SWE-Bench-Lite
($+3.67/+4.00$). Both Qwen2.5-Coder-Instruct sizes benefit from the same
mid-training corpus and recipe, indicating that the structural prior is
not absorbed by the larger pretrained checkpoint within this family.

\textbf{Transfer across post-training pipelines.}
Replacing R2E-Gym with SWE-Smith on the same 7B base yields $+5.30$
points on Verified, larger than the $+2.80$ under R2E-Gym; on Lite the
SWE-Smith pairing gains only $+0.50$, smaller than $+3.67$ for the
R2E-Gym pairing. The two pipelines together indicate that mid-training
is not tuned to a single post-training data distribution, though the
magnitude of its benefit depends on the baseline pipeline it composes
with.

\textbf{Transfer to a non-Qwen2.5 base.}
Switching to Qwen3-8B paired with SWE-Lego, mid-training improves
Verified by $+3.20$ and Lite by $+5.40$, comparable to the Qwen2.5-Coder
gains. This single comparison varies the post-training pipeline jointly
with the base model (SWE-Lego is required because Qwen3-8B's
$40{,}960$-token context exceeds the R2E-Gym/SWE-Smith scaffolds), so the
result should be read as ``not specific to the Qwen2.5-Coder $+$
R2E-Gym/SWE-Smith pairing'' rather than as a guarantee across
base-model families.

% =============================================================================
% 4.3 Capability preservation and cross-domain transfer
% =============================================================================
\vspace{-0.5ex}
\subsection{Capability Preservation and Cross-Domain Transfer}
\label{sec:exp:gen}

A natural concern with task-specialized post-training is that it erodes
capabilities the base model already had. We test this on six benchmarks
outside SWE-Bench and---for $\tau$-bench and BFCL---outside the coding
domain entirely. Because the controlled triple (instruct base vs.\
post-training only vs.\ mid-training $+$ post-training) is the most
expensive to run, we conduct it at 14B with R2E-Gym
(Table~\ref{tab:cap}).

% -----------------------------------------------------------------------------
% Capability preservation table
% -----------------------------------------------------------------------------
\begin{table}[t]
\centering
\caption{Capability preservation and cross-domain transfer at 14B with
R2E-Gym. Bold marks the better trained variant per column (post-only
vs.\ ours); the Instruct row is shown as a reference ceiling. All cells
are percentages.}
\label{tab:cap}
\small
\setlength{\tabcolsep}{4pt}
\begin{tabular}{lcccccc|c}
\toprule
& \multicolumn{3}{c}{Non-agent coding}
& \multicolumn{1}{c}{Agent OOD}
& \multicolumn{2}{c}{Tool use}
& \\
\cmidrule(lr){2-4} \cmidrule(lr){5-5} \cmidrule(lr){6-7}
Setting
& LiveCode & OJBench & FSB-EN
& Terminal
& $\tau$-bench & BFCL
& Avg \\
\midrule
Instruct Model
  & 37.20 & 5.20 & 53.80
  & 0.00
  & 5.70 & 23.20
  & 20.85 \\
+ R2E-Gym
  & 24.10 & 2.80 & 47.72
  & 2.41
  & 3.40 & 15.80
  & 16.04 \\
\textbf{+ FIM Mid-Train + R2E-Gym}
  & \textbf{35.20} & \textbf{4.74} & \textbf{48.25}
  & \textbf{3.66}
  & \textbf{7.30} & \textbf{18.20}
  & \textbf{19.56} \\
\rowcolor{blue!8}
$\Delta$ (vs.\ R2E-Gym only)
  & $+11.10$ & $+1.94$ & $+0.53$
  & $+1.25$
  & $+3.90$ & $+2.40$
  & $+3.52$ \\
\bottomrule
\end{tabular}
\end{table}

\textbf{Agentic post-training has a substantial hidden capability cost.}
R2E-Gym alone reduces every non-agent and tool-use benchmark relative to
the Instruct ceiling: LiveCodeBench by $13.10$, BFCL by $7.40$,
FullStackBench-EN by $6.08$, $\tau$-bench by $2.30$, and OJBench by
$2.40$. Averaged across the six benchmarks the post-trained model loses
$4.81$ points relative to Instruct---the implicit cost paid for
SWE-Bench gains, rarely highlighted in agent papers.

\textbf{Mid-training largely closes the regression gap.}
Adding FIM mid-training before the same post-training restores
LiveCodeBench by $+11.10$, OJBench by $+1.94$ (within $0.46$ of the
Instruct ceiling), and FullStackBench-EN by $+0.53$. The six-benchmark
average rises from $16.04$ to $19.56$ ($+3.52$ over post-training only),
while the SWE-Bench-Verified gain on the same base is preserved
(Table~\ref{tab:main}). Mid-training therefore improves the
cost--benefit profile of agentic post-training: in-domain target
metrics improve and the bulk of off-distribution erosion is undone in
the same training run.

\textbf{Cross-domain transfer to non-coding tool use.}
$\tau$-bench and BFCL contain no Python code-editing data, and our
mid-training corpus contains no tool-use trajectories. Mid-training
nevertheless improves $\tau$-bench by $+3.90$ and BFCL by $+2.40$ over
post-training alone, with a directionally consistent recovery on
Terminal-Bench ($+1.25$). Because the mid-training data carries no
tool-use signal, the only mechanism by which it can affect these
benchmarks is by installing a structural prior at mid-training time
that survives subsequent post-training---direct evidence for the
function-call/tool-call isomorphism (Section~\ref{sec:method:motivation}).

\vspace{-0.5ex}
\subsection{Ablation Studies}
\label{sec:exp:ablation}

We run two controlled ablations on Qwen2.5-Coder-7B-Instruct with R2E-Gym
post-training and a shared mid-training-free baseline. The first
isolates the role of the chain-of-thought rationale; the second isolates
the role of the function-aware selection pipeline. All selection
variants in block~(B) operate under a controlled budget---each variant
selects exactly $200$K single functions, with hard filters
(LoC~$\in[10,200]$, dunder methods excluded) applied uniformly---so
differences across rows reflect \emph{which} functions are selected
rather than \emph{how many}. Block~(C) varies mask granularity. Due to
compute constraints, ablations run at 7B only.

% -----------------------------------------------------------------------------
% Ablation table (unchanged)
% -----------------------------------------------------------------------------
\begin{table}[t]
\centering
\caption{Ablations on the 7B model with R2E-Gym post-training.
\textbf{(A)}~rationale source. \textbf{(B)}~function-selection algorithm.
\textbf{(C)}~mask granularity (single vs.\ multi-function groups).
All blocks share the baseline (\emph{w/o mid-train}) and a controlled
$200$K-target budget; CoT is fixed to Gemini-3-Flash in (B) and (C).
Bold marks the best configuration per block; the $80\%/15\%/5\%$ mixture
in~(C) is the recipe used in our main results (Table~\ref{tab:main}),
which is trained on the full corpus rather than the $200$K budget here.}
\label{tab:ablation}
\small
\setlength{\tabcolsep}{8pt}
\begin{tabular}{lccc}
\toprule
Setting & SWE-Bench-Verified & SWE-Bench-Lite & Average \\
\midrule
\multicolumn{4}{l}{\emph{(A) Chain-of-thought source (selection fixed to ``full'')}} \\
\midrule
w/o mid-train (baseline)                  & 15.00 & 11.33 & 13.17 \\
+ FIM, no CoT                             & 15.60 & 11.00 & 13.30 \\
+ FIM, self-CoT (Qwen2.5-Coder-7B-Instruct) & 16.40 & 13.30 & 14.85 \\
\textbf{+ FIM, Gemini-3-Flash CoT}        & \textbf{17.00} & \textbf{14.20} & \textbf{15.60} \\
\midrule
\multicolumn{4}{l}{\emph{(B) Function-selection algorithm (CoT fixed to Gemini-3-Flash)}} \\
\midrule
w/o mid-train (baseline)                    & 15.00 & 11.33 & 13.17 \\
Random                                      & 15.30 & 12.60 & 13.95 \\
Gemini-selected                             & 16.40 & 13.70 & 15.05 \\
PDG only                                    & 16.10 & 13.60 & 14.85 \\
PDG + complexity ($\hat{H}$)                & 16.50 & 13.60 & 15.05 \\
PDG + inferability ($\hat{I}$)              & 16.70 & 14.00 & 15.35 \\
\textbf{Full (PDG + $\hat{H}$ + $\hat{I}$)} & \textbf{17.00} & \textbf{14.20} & \textbf{15.60} \\
\midrule
\multicolumn{4}{l}{\emph{(C) Mask granularity (selection fixed to ``full'', CoT fixed to Gemini-3-Flash)}} \\
\midrule
w/o mid-train (baseline)             & 15.00 & 11.33 & 13.17 \\
Single-function only                 & 17.00 & 14.20 & 15.60 \\
85\% single + 15\% pair ($k=2$)       & 17.20 & 14.60 & 15.90 \\
95\% single + 5\% triple ($k=3$)      & 17.00 & 14.40 & 15.70 \\
\textbf{80\% single + 15\% pair + 5\% triple} & \textbf{17.40} & \textbf{14.80} & \textbf{16.10} \\
\bottomrule
\end{tabular}
\end{table}

\textbf{FIM structure contributes independently of CoT distillation.}
Block~(A) addresses a natural concern that improvements stem mainly
from distilling reasoning. Removing the rationale entirely (\emph{no
CoT}) lifts the average by only $+0.13$, with $+0.60$ on Verified---a
small effect consistent with the hypothesis that much of the underlying
code is already familiar to the base model from pretraining, leaving
little headroom for an unguided FIM signal alone. The more
informative comparison is against a frontier teacher.
Replacing Gemini-3 with rationales
generated by the model under training (\emph{self-CoT}) reaches
$14.85$, recovering $1.68$ of the $2.30$-point gap to \emph{Gemini-3
CoT}: high-quality external rationales help, but a sizable fraction of
the gain is reproducible without a frontier teacher. The function-aware
FIM objective is therefore effective even when paired with the model's
own rationales, and our recipe does not merely distill from a frontier
teacher.

\textbf{Function-selection algorithm is the dominant lever.}
Block~(B) varies the selection algorithm with CoT and budget held
fixed. \emph{Random} masking establishes the floor at $13.95$, a margin
of $+0.78$ over no mid-training: even random function-level masking on
top of the basic hard filters helps modestly. Letting Gemini-3 choose
which functions to mask (\emph{Gemini-selected}) reaches $15.05$
($+1.10$ over Random), showing that frontier judgment on \emph{which}
function to mask is helpful but not by itself sufficient. Restricting
the candidate pool to functions with at least one PDG neighbor
(\emph{PDG only}) reaches $14.85$, between Random and Gemini-selected:
structural connectivity alone contributes a small gain. Adding a
scoring function on top of the PDG pool yields the next layer:
\emph{PDG + $\hat{H}$} reaches $15.05$ and \emph{PDG + $\hat{I}$}
reaches $15.35$, showing that intrinsic difficulty and contextual
recoverability each contribute beyond the structural filter. Combining
both via the harmonic-mean-like form (\emph{Full}) reaches $15.60$
($+0.25$ over the better single-score variant), so $\hat{H}$ and
$\hat{I}$ measure complementary properties and their combination is not
redundant.

\textbf{Mask granularity: pair masking helps, triples saturate.}
Block~(C) mixes single-function targets with multi-function groups
(Section~\ref{sec:method:groups}). Adding $15\%$ pair targets on top of
single functions raises the average from $15.60$ to $15.90$ ($+0.30$);
substituting $5\%$ triple targets instead is essentially neutral
($15.70$). The full $80\%/15\%/5\%$ mix achieves $16.10$ ($+0.50$ over
single-only), with the largest gain on Lite ($+0.60$). The pattern is
consistent with the analysis in Section~\ref{sec:analysis:multifn}:
training on cross-function dependencies disproportionately helps tasks
whose gold patches span multiple functions, but the marginal return
from larger groups diminishes as coupling becomes harder to maintain in
the joint-masked context.
